\section{SSI for IoT Devices}

\subsubsection*{IoT Use Case and Decentralization in SSI}

This section describes how Self-Sovereign Identity (SSI) can be implemented in
IoT ecosystems, how embedded wallets operate inside constrained devices, and why
the SSI model is inherently decentralized. The focus is on identity anchoring,
device–issuer verification models, and implications for key management.

\paragraph{IoT Devices and Embedded Wallets}
In IoT systems, unlike user-centric web scenarios, no human interacts directly
with a wallet UI. \\
The SSI model must therefore be implemented in software running on the device,
as an \textit{embedded wallet}.  \\
This component handles:
\begin{itemize}
    \item key generation and storage,
    \item DID creation and DID Document publication,
    \item verifiable credential management (issuance, refresh, rotation),
    \item interaction with issuers and verifiers,
    \item secure signing of presentations.
\end{itemize}

\noindent
IoT deployments considered here operate in a \textbf{closed world} under a
single administrator, contrary to open web environments.\\
TLS is still used to create secure channels, typically with server-side X.509
authentication; SSI is used on top for client authentication at application
layer.


\paragraph{Combining TLS and SSI}
Even in IoT, the device must usually authenticate the server via TLS, using
traditional X.509 certificates.\\  
Client authentication can then occur through verifiable presentations.  
This results in a combination of two identity technologies:
\begin{itemize}
    \item TLS certificates for secure channel establishment,
    \item SSI credentials for application-layer device identity.
\end{itemize}

\noindent
However, in IoT scenarios the dependence on X.509 may be reduced, and SSI-based credentials can be used as certificate-like artifacts at lower layers of the stack. This reduces the overhead of certificate lifecycle management.

\paragraph{Anchoring Identity in Hardware}
When available, hardware secure elements such as TPMs or other Secure Elements (SE) can be used to generate and protect the device's key pair.\\ 
The private key remains inside the hardware; the corresponding public key is ncoded in the DID Document and published through the DID method.
\smallskip


\noindent
Using hardware strengthens:
\begin{itemize}
    \item key protection,
    \item integrity of cryptographic operations,
    \item resistance to device compromise.
\end{itemize}

\noindent
The embedded wallet must be able to use the hardware keys for all VC- and
VP-related signatures.

\paragraph{Verification Approaches for Issuing Credentials}
Before issuing a verifiable credential to an IoT device, the issuer must verify the device identity. Several approaches exist:
\bigskip

\noindent
\textbf{1. Remote Attestation (correct approach).}  
If a TPM is present, the issuer can perform remote attestation to verify:
\begin{itemize}
    \item device integrity,
    \item software configuration,
    \item the authenticity of the TPM,
    \item the origin and residency of the key pair inside the TPM.
\end{itemize}
Once attestation succeeds, the issuer may safely issue a verifiable credential bound to the DID created by the device.
\bigskip

\noindent
\textbf{2. Forcing key usage (incorrect approach).}  
An administrator may want to \emph{force} devices to use keys stored in the secure element by verifying, through TPM attestation keys, that a given key was generated and resides inside the TPM.\\
However, this creates a problem: the issuer must bind the VC to a DID method that prevents autonomous key rotation.
\smallskip

\noindent
Two options appear:
\begin{itemize}
    \item DLT-based DID methods,
    \item DID:key methods.
\end{itemize}

\noindent
DLT methods allow the device to later update its DID Document with a new public key—thus bypassing the enforced hardware key.\\
DID:key binds the credential directly to the current key, preventing rotation.
\medskip

\noindent
Both approaches violate SSI principles, because:
\begin{itemize}
    \item the issuer ends up controlling which key the device must use,
    \item the device loses the ability to rotate keys autonomously,
    \item decentralization is broken.
\end{itemize}

\noindent
Therefore, such enforcement is incompatible with SSI philosophy.

\paragraph{Correct Interpretation of SSI Principles}
SSI requires that:
\begin{itemize}
    \item the device controls the key pair,
    \item the device selects the DID method,
    \item issuers only validate identity facts, they do not mandate key usage,
    \item the holder remains free to rotate keys at any time.
\end{itemize}
\medskip

\noindent
Any issuer-side policy that forces a device to use one specific key (e.g., a TPM key) breaks the decentralized architecture.

\paragraph{Alternative Hardware Anchoring: PUFs}
Besides TPMs, identity can be anchored using \textbf{Physical Unclonable
Functions (PUFs)}.\\
A PUF extracts unique hardware characteristics of the board to derive keys that:
\begin{itemize}
    \item cannot be cloned,
    \item cannot be transplanted to another device,
    \item inherently bind identity to physical hardware.
\end{itemize}
\medskip

\noindent
This offers an alternative way to build strong IoT identities without storing long-term secrets in software.

\paragraph{People-to-People and People-to-Things Interactions}
Once SSI is implemented correctly, any entity—human or device—can interact with any other peer in the ecosystem.\\
Interoperability is guaranteed by the standardization of:
\begin{itemize}
    \item DID methods,
    \item Verifiable Credential data models,
    \item OpenID for Verifiable Credential Issuance (OIDC4VCI),
    \item OpenID for Verifiable Presentations (OIDC4VP).
\end{itemize}
\medskip

\noindent
Thus, verifiers and holders can be humans, servers, services, or embedded devices.

\paragraph{Why SSI Is Decentralized}
SSI is decentralized because:
\begin{itemize}
    \item No single authority controls all components of identity.
    \item The holder generates and manages its own key pair.
    \item The holder selects the DID method and controls the DID.
    \item Issuers only provide verifiable credentials, not the identity itself.
    \item No entity can fully revoke another peer's identity.
\end{itemize}

For example:
\begin{itemize}
    \item The issuer can revoke only the credential it issued.
    \item The holder can revoke its DID at any time.
    \item Key rotation is autonomous and does not require issuer intervention.
\end{itemize}
\smallskip

\noindent
This separation between:
\[
\text{public key (ledger)} \quad \text{and} \quad \text{identity attributes (VC)}
\]
is the fundamental design that enables decentralization.

\paragraph{Advantages for Large-Scale Systems}
The model simplifies identity lifecycle management:
\begin{itemize}
    \item Devices can rotate keys independently.
    \item Devices can revoke their DID instantly.
    \item No need for certificate revocation lists or CA-managed PKI.
    \item Administrators perform fewer manual identity maintenance operations.
\end{itemize}

\noindent
As a consequence, SSI reduces operational costs and scales better in large IoT deployments.



